\hypertarget{chapter-07-page-145}{}
\subsection{$A_p$ 条件}
\label{sec:ch7-ap-condition}

本章要找出所有非负局部可积函数 $w$, 使此前研究的算子在空间 $L^p(w)$ 上有界. 这些空间是把 $L^p$ 空间中的 Lebesgue 测度换成测度 $w\,dx$ 得到的. 我们称这样的函数为权; 对可测集 $E$, 定义 $w(E)=\int_Ew$.

首先假设 Hardy--Littlewood 极大函数满足带权弱型不等式, 并由此推出权 $w$ 的必要条件. 为简化分析, 本章把 Hardy--Littlewood 极大函数写成
\[
 Mf(x)=\sup_{Q\ni x}\frac1{|Q|}\int_Q|f(y)|\,dy,
\]
其中上确界取遍所有包含 $x$ 的立方体 $Q$. 这是 \eqref{eq:ch2-hl-maximal-uncentered-cube} 定义的非中心极大算子, 它与原定义逐点等价.

$M$ 关于 $w$ 的带权弱型 $(p,p)$ 不等式为
\begin{equation}
\label{eq:ch7-weighted-weak-maximal}
 w(\{x\in\mathbb R^n:Mf(x)>\lambda\})
 \le\frac C{\lambda^p}\int_{\mathbb R^n}|f(x)|^pw(x)\,dx.
\end{equation}
设 $f$ 非负, $Q$ 是满足 $f(Q)=\int_Qf>0$ 的立方体. 固定 $0<\lambda<f(Q)/|Q|$. 则 $Q\subset\{x:M(f\chi_Q)(x)>\lambda\}$, 所以由 \eqref{eq:ch7-weighted-weak-maximal},
\[
 w(Q)\le\frac C{\lambda^p}\int_Q|f(x)|^pw(x)\,dx.
\]

\hypertarget{chapter-07-page-146}{}
令 $\lambda\uparrow f(Q)/|Q|$, 得
\begin{equation}
\label{eq:ch7-testing-inequality}
 w(Q)\left(\frac{f(Q)}{|Q|}\right)^p
 \le C\int_Q|f|^pw.
\end{equation}
给定可测集 $S\subset Q$, 在 \eqref{eq:ch7-testing-inequality} 中取 $f=\chi_S$, 得
\begin{equation}
\label{eq:ch7-weight-set-comparison}
 w(Q)\left(\frac{|S|}{|Q|}\right)^p\le Cw(S).
\end{equation}
由此立即得到: 权 $w$ 或恒为零, 或几乎处处为正; 权 $w$ 或局部可积, 或几乎处处为无穷. 因而, 尽管我们预先假设 $w$ 局部可积, 这其实是加权范数不等式的推论.

先考虑 $p=1$. 此时 \eqref{eq:ch7-weight-set-comparison} 化为
\[
 \frac{w(Q)}{|Q|}\le C\frac{w(S)}{|S|}.
\]
令 $a=\operatorname*{ess\,inf}_{x\in Q}w(x)$. 对每个 $\varepsilon>0$, 存在具有正测度的 $S_\varepsilon\subset Q$, 使 $S_\varepsilon$ 上 $w(x)<a+\varepsilon$. 因此
\[
 \frac{w(Q)}{|Q|}\le C(a+\varepsilon).
\]
令 $\varepsilon\downarrow0$, 得对任意立方体 $Q$,
\begin{equation}
\label{eq:ch7-a1-cube-condition}
 \frac{w(Q)}{|Q|}\le Cw(x)
 \quad\text{对几乎每个 }x\in Q.
\end{equation}
这称为 $A_1$ 条件, 满足它的权称为 $A_1$ 权. 条件 \eqref{eq:ch7-a1-cube-condition} 等价于
\begin{equation}
\label{eq:ch7-a1-maximal-condition}
 Mw(x)\le Cw(x)
 \quad\text{对几乎每个 }x\in\mathbb R^n.
\end{equation}
显然 \eqref{eq:ch7-a1-maximal-condition} 蕴含 \eqref{eq:ch7-a1-cube-condition}. 反之, 若 \eqref{eq:ch7-a1-cube-condition} 成立, 则对有理顶点立方体取可数并可知 $Mw(x)>Cw(x)$ 只在零测集上成立.

\hypertarget{chapter-07-page-147}{}
现在设 $1<p<\infty$. 在 \eqref{eq:ch7-testing-inequality} 中取 $f=w^{1-p'}\chi_Q$, 得到等价条件
\begin{equation}
\label{eq:ch7-ap-condition}
 \left(\frac1{|Q|}\int_Qw\right)
 \left(\frac1{|Q|}\int_Qw^{1-p'}\right)^{p-1}\le C,
\end{equation}
其中 $C$ 与 $Q$ 无关. 推导时若尚不知道 $w^{1-p'}$ 局部可积, 可先用 $\min(w^{1-p'},N)$ 代替, 再令 $N\to\infty$. 由于 $w>0$ 几乎处处, \eqref{eq:ch7-ap-condition} 本身会推出 $w^{1-p'}$ 局部可积.

条件 \eqref{eq:ch7-ap-condition} 称为 $A_p$ 条件, 满足它的权称为 $A_p$ 权.

\begin{Theorem}
\label{thm:ch7-maximal-weighted-weak}
设 $1\le p<\infty$. 弱型 $(p,p)$ 不等式 \eqref{eq:ch7-weighted-weak-maximal} 成立, 当且仅当 $w\in A_p$.
\end{Theorem}

\begin{Proof}
必要性已在上面证明. 为证充分性, 先考虑 $p=1$. 定理~\ref{thm:ch2-weighted-maximal} 的带测度极大函数弱型估计右端含 $Mw$; 由于 $w\in A_1$, 可用 \eqref{eq:ch7-a1-maximal-condition} 把 $Mw$ 换成 $Cw$.

现在设 $p>1$ 且 $w\in A_p$. 先由 Hölder 不等式和 \eqref{eq:ch7-ap-condition} 得
\[
 \left(\frac1{|Q|}\int_Q|f|\right)^p
 \le\left(\frac1{|Q|}\int_Q|f|^pw\right)
 \left(\frac1{|Q|}\int_Qw^{1-p'}\right)^{p-1}
 \le\frac C{w(Q)}\int_Q|f|^pw,
\]
即 \eqref{eq:ch7-testing-inequality}, 从而也有 \eqref{eq:ch7-weight-set-comparison}.

固定非负 $f\in L^p(w)$, 并先设 $f\in L^1$. 在高度 $4^{-n}\lambda$ 对 $f$ 作 Calderón--Zygmund 分解, 得到两两不交的立方体 $\{Q_j\}$, 满足 $f(Q_j)>4^{-n}\lambda|Q_j|$. 与引理~\ref{lem:ch2-dyadic-to-cube-maximal} 的证明相同,
\[
 \{x\in\mathbb R^n:Mf(x)>\lambda\}\subset\bigcup_j3Q_j.
\]

\hypertarget{chapter-07-page-148}{}
因此
\[
\begin{aligned}
 w(\{Mf>\lambda\})
 &\le\sum_jw(3Q_j)
 \le C3^{np}\sum_jw(Q_j)\\
 &\le C3^{np}\sum_j\left(\frac{|Q_j|}{f(Q_j)}\right)^p
 \int_{Q_j}|f|^pw
 \le C3^{np}\left(\frac{4^n}\lambda\right)^p
 \int_{\mathbb R^n}|f|^pw.
\end{aligned}
\]
一般的 $f\in L^p(w)$ 可由截断得到, 常数与截断无关.
\end{Proof}

\begin{Proposition}
\label{prop:ch7-ap-properties}
有下列性质:
\begin{enumerate}[label=(\arabic*)]
\item $A_p\subset A_q$, $1\le p<q$;
\item $w\in A_p$ 当且仅当 $w^{1-p'}\in A_{p'}$;
\item 若 $w_0,w_1\in A_1$, 则 $w_0w_1^{1-p}\in A_p$.
\end{enumerate}
\end{Proposition}

\begin{Proof}
若 $p=1$, 则
\[
 \left(\frac1{|Q|}\int_Qw^{1-q'}\right)^{q-1}
 \le\operatorname*{ess\,sup}_{x\in Q}w(x)^{-1}
 \le C\left(\frac{w(Q)}{|Q|}\right)^{-1}.
\]
若 $p>1$, 第一条由 Hölder 不等式直接得到. 第二条中, $w^{1-p'}$ 的 $A_{p'}$ 条件正是 $w$ 的 $A_p$ 条件的 $p'-1$ 次幂.

为证第三条, 只需证明
\begin{equation}
\label{eq:ch7-a1-factorization-ap}
 \left(\frac1{|Q|}\int_Qw_0w_1^{1-p}\right)
 \left(\frac1{|Q|}\int_Qw_0^{1-p'}w_1\right)^{p-1}\le C.
\end{equation}

\hypertarget{chapter-07-page-149}{}
由 $A_1$ 条件, 对 $x\in Q$ 和 $i=0,1$,
\[
 w_i(x)^{-1}
 \le\operatorname*{ess\,sup}_{x\in Q}w_i(x)^{-1}
 =\left(\operatorname*{ess\,inf}_{x\in Q}w_i(x)\right)^{-1}
 \le C\left(\frac{w_i(Q)}{|Q|}\right)^{-1}.
\]
把它代入 \eqref{eq:ch7-a1-factorization-ap} 中具有负指数的因子, 即得所需不等式.
\end{Proof}
